\section{A charge for non-skew multiplication gates}\label{app:nonskew}

This appendix extends the polar count to a circuit model with a stated
budget on non-skew multiplication gates. A binary multiplication
gate is called \emph{non-skew} if neither child is an input variable or
field constant. Other multiplication gates are skew. Circuits are
commutative and division-free; arbitrary fan-out and cancellation are
allowed. The related noncommutative restrictions studied
in~\cite{LMS} are different models, and their lower bounds are not used.

\begin{theorem}[Circuit-state polar count]\label{thm:nonskew}
Let $f$ be a homogeneous degree-$d$ form in $k\ge3$ variables with
smooth projective hypersurface. Suppose $f$ is a fixed affine
restriction of a coefficient limit of outputs of circuits with $S$
internal binary gates, $N$ input variables, and at most $q$ non-skew
multiplication gates, where $0\le q\le S$. A larger stated budget may
always be replaced by $S$. Put $M=S+N+1$. Then
\[
 d(d-1)^{k-2}\le T_q(M,k),
\]
where $T_0(M,k)=B(M,k)$ and, for $q\ge1$,
\begin{equation}\label{eq:nonskew-count}
 T_q(M,k)\le
 2^{q+k-2}\binom{q+k-3}{k-2}\binom{2M-1}{k-1}.
\end{equation}
In particular, for every fixed $C\ge0$, the budget $q\le C(k-1)$
implies $M=\Omega_C(kd)$ as $d$ tends to infinity.
\end{theorem}

\begin{corollary}\label{cor:nonskew-per}
For each fixed constant $C$, a division-free circuit computing
$\per_n$ with at most $Cn^2$ non-skew multiplication gates has
$\Omega_C(n^2\log n)$ internal gates. The same holds for coefficient
approximation with these size and non-skew-gate budgets.
\end{corollary}

\begin{proof}
Use the fixed smooth restriction from the proof of
the first part of Theorem~\ref{thm:main}, taking
$t=\lfloor\log_2 n\rfloor$, $d=\lfloor t/2\rfloor$,
$k=\lfloor n^2/16\rfloor$ for all sufficiently large $n$,
and $N=n^2$ original
input registers. Affine restriction changes their labels to affine
forms but preserves their role as skew operands in the construction
below. Theorem~\ref{thm:nonskew} gives
$S+N+1=\Omega_C(n^2\log n)$; subtracting $N+1$ proves the assertion
for all sufficiently large $n$.
\end{proof}

\subsection*{State equations and their adjoints}

Introduce one state $w$ for each original input variable and each
internal gate, for a total of $M-1$ states. After restriction, the input
equations are $w_i=L_i(x)$ with $L_i$ affine. In topological order, each
addition, subtraction, or skew multiplication has an equation affine
in the states, with coefficients affine in $x$. For a skew product,
replace its input child by the affine label $L_i(x)$; retain the state
register for the other operand. Constants require no state register.
Each non-skew product has an equation
\[
 w_g-w_aw_b=0.
\]
The Jacobian of these $M-1$ graph equations with respect to the states
is triangular with diagonal one. Thus they define the graph of the
computation, including every use of every shared value.

Add the output equation $w_{\rm out}=0$. On
$\PP^k_{[t:x]}\times\PP^{M-1}_{[v_0:w]}$, the $M-q$ linear-state
equations, including the output, become forms $F_i$ of bidegree
$(1,1)$. Examples are
\[
 tv_i-\widehat L_i(x)v_0,\qquad
 tv_g-\widehat L_i(x)v_a,\qquad tv_{\rm out}.
\]
For the other $q$ equations use
\[
 tQ_j(v)=t(v_gv_0-v_av_b),
\]
of bidegree $(1,2)$. The added factors $t$ have no effect on the
chart $t=v_0=1$. Their extra components at infinity are permitted.

At a zero of the output with nonzero differential, the Jacobian of
all $M$ constraints in the $M-1$ state directions has rank $M-1$.
Its left kernel is a unique multiplier line. Write the multipliers as
$u_i$ for the $F_i$ and $\lambda_j$ for the $tQ_j$. The multiplier of
the output equation is nonzero: otherwise the invertible graph
Jacobian forces every multiplier to vanish. Normalize that multiplier
to one. Implicit differentiation of the graph equations shows that
\[
 S_y=\sum_i u_i(F_i)_y+\sum_j\lambda_j(tQ_j)_y
\]
vanishes for all state directions, and its $x$-components are, up to
sign, the differential of the output. The state adjoints sum all
contributions from fan-out; no parse-tree expansion occurs.

Choose a simple polar fiber as in Proposition~\ref{prop:polar}. It
has $\delta=d(d-1)^{k-2}$ points. At each point the graph equations
eliminate the states and the adjoint equations eliminate the normalized
multipliers, by invertible Jacobians. The resulting local system is
the original square polar system. Every lifted point is therefore
isolated and reduced. The same is true for the $\delta$ points
persisting in any sufficiently close polynomial circuit output.

\subsection*{A bundle count that respects reuse}

Write $H,V$ for the hyperplane classes of
$Y=\PP^k\times\PP^{M-1}$ and use the quotient projective bundle
\[
 \pi:\PP(E')\longrightarrow Y,\qquad
 E'=\mathcal O_Y^{\oplus(M-q)}\oplus
       \mathcal O_Y(V)^{\oplus q},\qquad
 L=c_1(\mathcal O_{\PP(E')}(1)).
\]
The coordinates $u_i$ have class $L$, whereas the $\lambda_j$ have
class $L-V$. Hence a state-adjoint equation has class $L+H$, and an
$x$-directional adjoint equation has class $L+V$. The latter has no
quadratic-constraint contribution, since $Q_j$ is independent of $x$;
it is nevertheless a section of the stated line bundle. The
homogenizing coordinate $t$ is not one of these $x$ directions.

The incidence has $M-q$ equations of class $H+V$, $q$ of class
$H+2V$, $M-1$ state adjoints of class $L+H$, $k-2$ polar adjoints of
class $L+V$, and one equation $\ell(x)-t$ of class $H$. Their number
is $k+2M-2=\dim\PP(E')$. All these bundles are globally generated:
$E'$ is globally generated, as are its twists by $H$ and $V$.
General perturbations of the sections thus have proper
zero-dimensional intersection. The isolated points just constructed
persist, giving the upper bound
\begin{equation}\label{eq:state-intersection}
 T_q(M,k)=\int_{\PP(E')}
 H(H+V)^{M-q}(H+2V)^q(L+H)^{M-1}(L+V)^{k-2}.
\end{equation}
This bound concerns the original system's isolated points and allows
excess components elsewhere. In particular, it does not assume
properness of the incidence at infinity.

For $q=0$ the bundle is trivial and
\eqref{eq:state-intersection} is $B(M,k)$ after exchanging the two
kernel-factor names. For $q\ge1$, the projective-bundle formula gives
\[
 \pi_*(L^{M-1+a})=\binom{q+a-1}{a}V^a\quad(a\ge0),\qquad
 \pi_*(L^j)=0\quad(j<M-1).
\]
For precision about conventions, these formulas follow from
$L^{M-q}(L-V)^q=0$ and $\pi_*(L^{M-1})=1$ by induction;
see~\cite[Sections 27.21 and 42.36]{Stacks}. No sign convention for
Segre classes is implicit.

All coefficients in the integrand and pushforward are nonnegative.
Bound $(H+2V)^q$ coefficientwise by $2^q(H+V)^q$. If $i$ copies
of $H$ are chosen from $(H+V)^M$ and $c$ from $(L+H)^{M-1}$, the
required $H$ degree is equivalent to $i+c=k-1$. Choosing $j$ copies
of $L$ from $(L+V)^{k-2}$ leaves a pushforward contribution bounded by
\[
 \sum_{j=c}^{k-2}\binom{k-2}{j}\binom{q+j-c-1}{j-c}
 \le 2^{k-2}\binom{q+k-3}{k-2}.
\]
The remaining sum is
\[
 \sum_{i+c=k-1}\binom Mi\binom{M-1}c
 =\binom{2M-1}{k-1},
\]
which proves~\eqref{eq:nonskew-count}. The count is nondecreasing
in $q$: both the constraint coefficients and the pushforward
coefficients increase. An upper budget on $q$ therefore suffices.

Fixed-size circuit families form a finite union of constructible
polynomial images. The simple target polar points persist under
coefficient approximation, with no boundedness assumption on circuit
parameters or state values. This proves the border assertion and
completes the proof of Theorem~\ref{thm:nonskew}.

\subsection*{An explicit tradeoff and its limit}

Taking $(k-1)$-st roots and using the standard binomial estimates
gives the convenient bound, for all $q\ge0$,
\begin{equation}\label{eq:state-tradeoff}
 M\ge\frac12+
 \frac{(k-1)(d-1)}
 {4e^2\,2^{q/(k-1)}(1+q/(k-2))}.
\end{equation}
For $q=0$, Proposition~\ref{prop:polar} gives a stronger constant.
For $q\ge1$, use
\[
 \binom{q+k-3}{k-2}^{1/(k-1)}
 \le e\left(1+\frac q{k-2}\right),\qquad
 \binom{2M-1}{k-1}^{1/(k-1)}
 \le\frac{e(2M-1)}{k-1}.
\]
If $q\le C(k-1)$, the denominator of
\eqref{eq:state-tradeoff} is bounded by a constant depending only on
$C$, proving the last assertion of the theorem.

This charge tolerates $O(k)$ non-skew gates but weakens as their number
grows. For example, repeated squaring computes
$\sum_{i=1}^k x_i^{2^a}$ using $k(a-1)$ non-skew squarings after the
first squares. Here $q/k$ grows with $\log d$, outside the fixed-$C$
regime. Thus the argument respects the earlier counterexample to a
linear charge of smooth-slice score to unrestricted circuit size.
Corollary~\ref{cor:nonskew-per} is a restricted circuit result, not an
unrestricted $\Omega(n^2\log n)$ circuit bound.

The accompanying finite checker independently expands 416 intersection
coefficients, checks 364 monic projective-bundle reductions and 32
repeated-squaring controls, and verifies 35 full local incidence
Jacobians, including 77 multiply-used internal states. These checks
support the displayed identities and local elimination; the preceding
argument proves the uniform statement.
